%!TEX TS-program = xelatex
%!TEX encoding = UTF-8 Unicode
% Built with Awesome-CV (https://github.com/posquit0/Awesome-CV)

\documentclass[11pt, a4paper]{awesome-cv}

\geometry{left=1.4cm, top=.8cm, right=1.4cm, bottom=1.8cm, footskip=.5cm}

\usepackage{xeCJK}
\setCJKmainfont{Hiragino Kaku Gothic ProN}
\setCJKsansfont{Hiragino Kaku Gothic ProN}
\setCJKmonofont{Hiragino Kaku Gothic ProN}

\colorlet{awesome}{awesome-red}
\setbool{acvSectionColorHighlight}{true}
\renewcommand{\acvHeaderSocialSep}{\quad\textbar\quad}

\name{Tomas}{Sierra Sanchez}
\position{Estudiante de Ingenier\'ia de Sistemas y Computaci\'on}
\address{Bogot\'a D.C, Colombia}

\mobile{+57 311 775 5877}
\email{t.sierras@uniandes.edu.co}

\begin{document}

\makecvheader[C]

\makecvfooter
  {\today}
  {Tomas Sierra Sanchez~~~\textperiodcentered~~~Hoja de Vida}
  {\thepage}

\cvsection{Educaci\'on}

\begin{cventries}

  \cventry
    {Ingenier\'ia de Sistemas y Computaci\'on, Lengua y Cultura Japonesa}
    {Universidad de los Andes}
    {Bogot\'a D.C, Colombia}
    {Fecha de graduaci\'on: Abril 2027}
    {}

  \cventry
    {Programa de Intercambio, Tsukuba Trans Pacific Program}
    {Universidad de Tsukuba}
    {Tsukuba, Jap\'on}
    {Septiembre 2025 -- Febrero 2026}
    {
      \begin{cvitems}
        \item {Viv\'i de primera mano la lengua y cultura japonesa, adem\'as de estudiar \'areas cercanas y relevantes a mi carrera en casa.}
      \end{cvitems}
    }

\end{cventries}

\cvsection{Experiencia}

\begin{cventries}

  \cventry
    {Desarrollador de Software -- Asistente de Investigaci\'on de Pregrado}
    {Nodo EdTech -- Universidad de los Andes}
    {Bogot\'a D.C, Colombia}
    {Abril 2026 -- Presente}
    {
      \begin{cvitems}
        \item {Contribuyo al desarrollo full-stack de la plataforma web No Estas Solo, trabajando tanto en el backend como en el frontend.}
        \item {Ayud\'e a mantener servicios en contenedores para apoyar el desarrollo y despliegue de la plataforma.}
      \end{cvitems}
    }

  \cventry
    {Tutor de Mastering Machine Learning}
    {Universidad de los Andes}
    {Bogot\'a D.C, Colombia}
    {Junio 2026 -- Julio 2026}
    {
      \begin{cvitems}
        \item {Ayud\'e con la planeaci\'on y el desarrollo de talleres para la clase.}
        \item {Apoy\'e a los estudiantes respondiendo inquietudes y calificando talleres y proyectos.}
      \end{cvitems}
    }

  \cventry
    {Asistente de Investigaci\'on de Pregrado en Inteligencia Artificial Generativa}
    {Facultad de Ingenier\'ia -- Universidad de los Andes}
    {Bogot\'a D.C, Colombia}
    {Enero 2025 -- Septiembre 2025}
    {
      \begin{cvitems}
        \item {Ayud\'e a desarrollar herramientas de IA generativa para mejorar la experiencia de los estudiantes en distintas clases de la facultad de ingenier\'ia, usando herramientas como N8N Automation y OpenWebUI.}
        \item {Realic\'e investigaci\'on a profundidad sobre ingenier\'ia de prompts para mejorar la eficiencia de estas herramientas.}
        \item {Colabor\'e con profesores y equipos de marketing de la universidad para implementar y probar estas herramientas en escenarios reales.}
      \end{cvitems}
    }

  \cventry
    {Monitor de Matem\'aticas Discretas}
    {Universidad de los Andes}
    {Bogot\'a D.C, Colombia}
    {Enero 2024 -- Junio 2025}
    {
      \begin{cvitems}
        \item {Brind\'e apoyo durante ejercicios pr\'acticos con estudiantes y ofrec\'i acompa\~namiento en la preparaci\'on de ex\'amenes mediante mentor\'ias y resoluci\'on de dudas.}
      \end{cvitems}
    }

  \cventry
    {Monitor de An\'alisis y Dise\~no de Algoritmos}
    {Universidad de los Andes}
    {Bogot\'a D.C, Colombia}
    {Agosto 2024 -- Diciembre 2024}
    {
      \begin{cvitems}
        \item {Asist\'i en ejercicios pr\'acticos con estudiantes y calific\'e tareas y quices de la clase.}
      \end{cvitems}
    }

  \cventry
    {Monitor de Introducci\'on a la Programaci\'on}
    {Organizaci\'on}
    {Bogot\'a D.C, Colombia}
    {Agosto 2023 -- Diciembre 2023}
    {
      \begin{cvitems}
        \item {Apoy\'e al profesor durante los talleres pr\'acticos resolviendo dudas de los estudiantes y calificando proyectos de clase.}
      \end{cvitems}
    }

\end{cventries}

\cvsection{Habilidades e Intereses}

\begin{cvskills}

  \cvskill
    {Lenguajes de Programaci\'on}
    {Python, Java, C++, C\#, TypeScript, JavaScript, Swift, Dart}

  \cvskill
    {Frameworks y Herramientas}
    {Next, Coolify, SQL, MongoDB, SpringBoot, Django, Flutter, XCode}

  \cvskill
    {Idiomas}
    {Espa\~nol (Lengua materna), Ingl\'es (IELTS 8.0 / C1), Franc\'es (DELF B1), Japon\'es (Nivel intermedio)}

\end{cvskills}

\end{document}
